% ===================================================================
%  TABLICA Nr 1 — Szkoła. — Liceum. — Klasa.
%
%  Ceci n'est PAS une transcription : c'est une traduction, faite en
%  2026, en polonais d'aujourd'hui. D'ou trois differences avec
%  texto/io et texto/fr, et trois seulement :
%
%   * pas de \nl ni de \cc — la traduction ne suit aucune ligne
%     imprimee, et les lignes de ce fichier ne sont que du confort ;
%   * pas de \begin{VUpage} — elle n'a ni feuillet ni folio, et la
%     page de lecture ne lui en affiche pas ;
%   * le gras se pose sur le TERME seul, ponctuation dehors.
%
%  Tout le reste est identique, et doit l'etre : les cles « %%K » sont
%  celles de texto/io, macro pour macro pour l'apparat, et les renvois
%  \textsuperscript{(n)} visent les MEMES objets de la planche — ce
%  sont eux qui ouvrent les gros plans.
%
%  UNE COLONNE HORS DES DEUX LISTES. Le polonais n'est ni dans les
%  vingt-neuf langues d'Ethnologue — il vient loin derriere le
%  bhojpouri — ni parmi les dix-sept pays de l'ido. Il est la parce
%  qu'on l'a demande, et texto/lingui.json le dit.
%
%  LA DEFENSE DE CETTE COLONNE SE JOUE A LA LETTRE, comme celle du
%  haoussa. Aucune voisine ne lui ressemble assez pour le contaminer :
%  le tcheque est ouest-slave comme lui, mais ě, ř, ů, č, š, ž ne sont
%  pas ą, ę, ł, ń, ó, ś, ź, ż, et un mot tcheque se verrait. Le danger
%  est le CLAVIER : neuf lettres ne s'obtiennent pas sans combinaison,
%  et une main pressee ecrit ksiazka, stol, dzien, reka la ou il faut
%  książka, stół, dzień, ręka. La faute ne se voit pas — le mot reste
%  lisible, il est seulement faux.
%
%  ET LE POLONAIS DECLINE, ce que ni l'ido ni le francais ne font. Le
%  renvoi appartient a son substantif et le suit dans son cas :
%  « na \VUgras{ławce} \textsuperscript{(32)} », locatif, et non au
%  nominatif parce que la planche nomme l'objet au nominatif. C'est le
%  mot qui porte le renvoi, pas le lemme.
%
%  L'ORDRE DES RENVOIS EST CELUI DE L'IDO, et le polonais le tient
%  sans peine : il pose son adjectif AVANT le nom comme l'ido, et sa
%  syntaxe est assez libre pour suivre. Les inversions du genitif —
%  « lekcja nauczyciela » — ne separent pas deux renvois voisins dans
%  ce tableau ; outils/parigi.py dit ou elles le feront.
%
%  LES PRENOMS SONT POLONISES, comme dans les autres colonnes :
%  Ioannes est Jan, Iakobus est Jakub, Karolus est Karol. Les noms
%  latins de la note (*) ne bougent pas : c'est d'eux que la note
%  parle.
% ===================================================================

%%K t01-apar-0 apar
\VUsaut{2.0mm}
\VUcentre{12.6pt}{120}{KSIĄŻECZKA OBJAŚNIAJĄCA}
\VUsaut{3.2mm}
\VUcentre{8.4pt}{160}{DO}
\VUsaut{2.6mm}
\VUtitre{15.6pt}{0}{tablic pomocniczych Delmas}
\VUsaut{3.4mm}
\VUornamento{9pt}
\VUsaut{5.0mm}
\VUcentre{13.2pt}{90}{SERIA PIERWSZA}
\VUsaut{3.0mm}
\VUfilet{16mm}
\VUsaut{5.6mm}

%%K t01-apar-1 apar
\VUtitre{15.0pt}{60}{TABLICA Nr 1}
\VUsaut{1.4mm}
\VUtitre{11.4pt}{0}{Szkoła, liceum, klasa.}
\VUsaut{2.2mm}
\VUfilet{20mm}
\VUsaut{6.4mm}

%%K t01-tit-1 sub
\VUpk{10.2pt}{40}{Opis ogólny.}
\VUsaut{3.0mm}

%%K t01-01-1 p
1. --- Ta tablica przedstawia \VUgras{klasę} w \VUgras{liceum}. W tej klasie stoi
osiem \VUgras{stołów} \textsuperscript{(18)} i osiem \VUgras{ławek} \textsuperscript{(32)}.
Na każdej ławce siedzi trzech uczniów. \VUgras{Nauczyciel} \textsuperscript{(7)} siedzi na
\VUgras{krześle} \textsuperscript{(8)} przy swojej \VUgras{katedrze} \textsuperscript{(6)}.

%%K t01-02-1 p
2. --- Na lewo od katedry stoi \VUgras{szafa} \textsuperscript{(15)}
\VUgras{oszklona}. Na prawo jest \VUgras{czarna tablica} \textsuperscript{(1)} z
\VUgras{gąbką} \textsuperscript{(2)} i \VUgras{kredą} \textsuperscript{(3)}.

%%K t01-02-2 p
Nad katedrą wiszą \VUgras{tablice ścienne} \textsuperscript{(9, 11, 12)} i
\VUgras{mapa} \textsuperscript{(10)}.

%%K t01-03-1 p
3. --- Nie wszyscy uczniowie siedzą na swoich \VUgras{miejscach} \textsuperscript{(17)}.
Jan \textsuperscript{(4)} stoi przy tablicy, trzyma gąbkę i ściera
\VUgras{figury geometryczne}.

%%K t01-03-2 p
Piotr jest blisko katedry, zbiera \VUgras{prace pisemne} \textsuperscript{(5)} i
zanosi je nauczycielowi.

%%K t01-04-1 p
4. --- Ten nauczyciel jest nauczycielem \VUgras{języków żywych}. Uczy uczniów
mówić, czytać i pisać w językach obcych \VUgras{(po niemiecku, angielsku,
hiszpańsku, włosku, rosyjsku} lub \VUgras{francusku)}.

%%K t01-05-1 p
5. --- Klasa jest bardzo obszerna i bardzo jasna. W głębi jest szerokie
\VUgras{okno} z dwoma \VUgras{lufcikami} \textsuperscript{(78)} i
\VUgras{oszklone drzwi}, żeby ją wietrzyć i oświetlać.

%%K t01-05-2 p
Ta klasa nie jest zimna. Pośrodku stoi, żeby ją ogrzewać, bardzo dobry
\VUgras{piec} \textsuperscript{(46)} ze swoją \VUgras{rurą} \textsuperscript{(45)}.

%%K t01-05-3 p
Do ściany przymocowane są \VUgras{wieszaki} \textsuperscript{(58)}, wiszą na nich
czapki i płaszcze uczniów.

%%K t01-tit-2 sub
\VUsaut{5.2mm}
\VUpk{10.2pt}{40}{Boisko}
\VUsaut{3.6mm}

%%K t01-06-1 p
6. --- \VUgras{Zegar} \textsuperscript{(63)} wskazuje piątą po dziewiątej.

%%K t01-06-2 p
\VUgras{Przerwa} \textsuperscript{(76)} właśnie się skończyła i lekcja zaczyna się
na nowo. Miejsce zabaw znajduje się na \VUgras{dziedzińcu} \textsuperscript{(75)}.
Widzimy kilku uczniów, którzy się bawią. \VUgras{Pomocnik nauczyciela}
\textsuperscript{(74)} ich pilnuje.

%%K t01-07-1 p
7. --- Na dziedzińcu widzimy ponadto różne osoby, mianowicie
\VUgras{inspektora} \textsuperscript{(73)}, dyrektora szkoły \VUgras{(rektora)}
\textsuperscript{(71)} i \VUgras{cenzora} \textsuperscript{(72)}.

%%K t01-07-2 p
Inspektor wizytuje szkoły, rektor kieruje liceum, cenzor czuwa nad naukami.

%%K t01-07-3 p
Czwartą osobą jest \VUgras{portier} \textsuperscript{(77)}. Niesie pod pachą rejestr,
żeby wpisywać nieobecnych.

%%K t01-08-1 p
8. --- W głębi dziedzińca są \VUgras{przyrządy gimnastyczne} \textsuperscript{(70)} i
warsztat do \VUgras{robót ręcznych} \textsuperscript{(69)}.

%%K t01-08-2 p
Na prawo od tablicy zaczynamy widzieć \VUgras{sale} do \VUgras{muzyki} i
\VUgras{rysunku}.

%%K t01-noto-1 noto
\VUnotes{143.2mm}{%
(*) Dla \textit{imion chrzestnych} radzi się przepisywać je również według formy
łacińskiej. To jedyny sposób, żeby uniknąć niezliczonych odmian. Jak zresztą wybierać
między \textit{Giovanni, Jean,} \textit{Johann, Jan, Hans, John, Ivan,} itd.? Czy
stworzymy osobny słownik imion chrzestnych? Czerpmy z łaciny ich mianownik i mówmy:
\VUgras{Ioannes, Ioanna; Iulius, Iulia; Iakobus; Andreas; Lukas.} (Gram. Zupełna).%
}

%%K t01-tit-3 sub
\VUpk{10.2pt}{40}{Opis szczegółowy.}
\VUsaut{2.4mm}

%%K t01-tit-4 sub
\VUpk{10.2pt}{40}{Dobrzy uczniowie.}
\VUsaut{3.0mm}

%%K t01-09-1 p
9. --- Uczniowie z pierwszej ławki na prawo od katedry to \VUgras{Henryk} (*),
\VUgras{Szczepan} i \VUgras{Karol}.

%%K t01-09-2 p
Henryk jest bardzo uważny i pracowity, słucha nauczyciela.

Na swoim stole ma \VUgras{zeszyt} \textsuperscript{(28)}, \VUgras{książkę}
\textsuperscript{(29)}, \VUgras{słownik} \textsuperscript{(30)} i \VUgras{piórnik}
\textsuperscript{(31)}. Jego książka ma wiele \VUgras{stron}, każdy rozdział zawiera
kilka \VUgras{ustępów}, a każdy \VUgras{wiersz} wiele \VUgras{wyrazów}.

%%K t01-09-4 p
Jego sąsiad Szczepan \textsuperscript{(24)} jest gorliwy. Zapisuje w swoim zeszycie
lekcję na następny raz. Ma piękny \VUgras{charakter pisma}. Jego książka jest
zamknięta. Widzimy tylko \VUgras{okładkę} \textsuperscript{(27)}. Obok niego leży jego
\VUgras{cyrkiel} \textsuperscript{(26)}.

%%K t01-09-5 p
Karol \textsuperscript{(23)} trzyma w ręku swoją książkę, otwiera ją, żeby powtórzyć
lekcję. Ma, jak każdy uczeń, \VUgras{kałamarz} \textsuperscript{(25)} przed sobą. Jest
posłuszny.

%%K t01-10-1 p
10. --- Przy sąsiednim stole są \VUgras{Ludwik}, \VUgras{Antoni} i \VUgras{Paweł}.
Ludwik powtarza swoją \VUgras{lekcję} \textsuperscript{(22)} i odpowiada na
\VUgras{pytania} nauczyciela. Antoni \textsuperscript{(21)} jest bardzo nieuważny, nawet
czasem nauczyciel go karze. Paweł \textsuperscript{(20)} jest dobrym kolegą, uprzejmym i
usłużnym. Jest bardzo uważny.

%%K t01-tit-5 sub
\VUsaut{4.2mm}
\VUpk{10.2pt}{40}{Źli uczniowie.}
\VUsaut{3.2mm}

%%K t01-11-1 p
11. --- Za Henrykiem siedzi \VUgras{Jerzy} \textsuperscript{(38)}. Nie jest złym
chłopcem, ale jest \VUgras{gadatliwy}, nieuważny i swawolny. Jego \VUgras{lekkomyślność}
i \VUgras{nieposłuszeństwo} sprawiają \VUgras{zamęt} podczas lekcji i często zasługuje
na surowe \VUgras{upomnienie}. Jego \VUgras{brulion} \textsuperscript{(35)} jest brudny,
\VUgras{plami} go atramentem swoim \VUgras{piórem} \textsuperscript{(34)}, dlatego ciągle
używa swojej \VUgras{gumki} \textsuperscript{(36)}; jego \VUgras{teczka}
\textsuperscript{(33)} jest podarta, bo jest bardzo niedbały. Po co przynosi swój
\VUgras{piórnik} \textsuperscript{(37)} do klasy języków żywych?

%%K t01-11-2 p
Jego sąsiad Edmund \textsuperscript{(39)} go nie lubi. Jest wprawdzie trochę gnuśny, ale
jest milczący i spokojny. Nie gada; ale zamiast słuchać, woli czytać
\VUgras{opowiadania} ze swojej \VUgras{czytanki}.

%%K t01-11-3 p
Trzecim uczniem z tej ławki jest \VUgras{Ludwik} \textsuperscript{(40)}. Jest pilny.
Pisze na białym \VUgras{arkuszu papieru}. Przed sobą położył swoją
\VUgras{bibułę} \textsuperscript{(41)}.

%%K t01-12-1 p
12. --- Uczniowie z drugiej ławki, na lewo od nauczyciela, są bardzo młodzi. Pierwszy,
mały \VUgras{Emil}, nie umie jeszcze bardzo dobrze pisać, przynosi swoją
\VUgras{tabliczkę} \textsuperscript{(42)}, swój \VUgras{rysik} i swoją \VUgras{gąbkę}
\textsuperscript{(43)}.

%%K t01-12-2 p
Obok niego są \VUgras{August} i \VUgras{Bernard} \textsuperscript{(44)}. Ten pracuje
dobrze, ale jego \VUgras{ubranie} jest bardzo zaniedbane.

%%K t01-13-1 p
13. --- Za nimi siedzą trzej \VUgras{leniwi}. \VUgras{Albert} nie uczy się swoich
lekcji. \VUgras{Jakub} \textsuperscript{(47)} śpi zamiast słuchać. Jego
\VUgras{lenistwo} ściąga na niego \VUgras{wymówki}, a nawet \VUgras{kary}. Wreszcie
\VUgras{Krystian} \textsuperscript{(48)} jest zbyt niepoważny. Źle się \VUgras{sprawuje}
i wcale nie stara się poprawić.

%%K t01-14-1 p
14. --- Jego sąsiedzi, przeciwnie, są dobrze wychowani. \VUgras{Ernest}
\textsuperscript{(51)} lubi porządek i regularność, zawsze używa swojej
\VUgras{linijki} \textsuperscript{(50)} i swojego \VUgras{ołówka} \textsuperscript{(49)}.
Jest \VUgras{eksternistą}.

%%K t01-14-2 p
\VUgras{Franciszek} \textsuperscript{(52)} jest \VUgras{internistą}, jest ubrany w
mundurek, je i sypia w liceum. Jest dobrym \VUgras{kolegą}.

%%K t01-14-3 p
Maurycy struga swój \VUgras{ołówek} swoim \VUgras{scyzorykiem} \textsuperscript{(56)},
jest bardzo staranny.

%%K t01-14-4 p
Przynosi wszystkie swoje książki, swoją \VUgras{gramatykę} \textsuperscript{(55)}, swój
\VUgras{notesik} \textsuperscript{(54)}, a nawet \VUgras{podkładkę do pisania}
\textsuperscript{(53)}. Jego \VUgras{tornister} \textsuperscript{(57)} stoi na podłodze
za nim.

%%K t01-14-5 p
Dwa stoły w głębi klasy zajmują \VUgras{dobrzy uczniowie} \textsuperscript{(19)}.

%%K t01-tit-6 sub
\VUsaut{4.6mm}
\VUpk{10.2pt}{40}{Rysunek i muzyka.}
\VUsaut{3.4mm}

%%K t01-15-1 p
15. --- W \VUgras{sali rysunkowej} są wdzięczne \VUgras{modele} zawieszone na ścianie i
\VUgras{lampa} \textsuperscript{(62)} z \VUgras{abażurem}, zawieszona u sufitu.
\VUgras{Nauczyciel} \textsuperscript{(60)} jest bardzo cierpliwy, poprawia
\VUgras{rysunki} \textsuperscript{(61)} uczniów i uczy ich rysować \VUgras{popiersie}
\textsuperscript{(59)} z natury.

%%K t01-16-1 p
16. --- \VUgras{Sala muzyczna} wydaje się bardzo ciekawa. Uczy się w niej czytać
\VUgras{nuty}, solfeżować \VUgras{dźwięki gamy} \textsuperscript{(67)} zapisane na
\VUgras{pięciolinii} (do, re, mi, fa, sol, la, si\,=\,C. D. E. F. G. A. B.), poznawać
\VUgras{klucze} i śpiewać urocze \VUgras{melodie}. Nuty kładzie się na \VUgras{pulpicie}
\textsuperscript{(65)}. Jeden z uczniów \VUgras{gra na fortepianie} \textsuperscript{(64)},
a \VUgras{nauczyciel muzyki} \textsuperscript{(66)} \VUgras{wybija takt}
\textsuperscript{(68)}.

%%K t01-17-1 p
17. --- Zaczynamy widzieć kąt \VUgras{warsztatu} do \VUgras{robót ręcznych}
\textsuperscript{(69)}, gdzie chłopcy mogą się uczyć heblować drewno i piłować żelazo.
Tego nie ma we wszystkich liceach.

%%K t01-tit-7 sub
\VUsaut{5.0mm}
\VUpk{10.2pt}{40}{Sala nauk ścisłych.}
\VUsaut{3.6mm}

%%K t01-18-1 p
18. --- Nauczyciel matematyki uczy uczniów \VUgras{geometrii, arytmetyki, rachunków} i
\VUgras{systemu metrycznego}. Widzimy na tablicy główne figury geometryczne:
\VUgras{koło} \textsuperscript{(a)}, \VUgras{kwadrat} \textsuperscript{(b)},
\VUgras{trójkąt} \textsuperscript{(c)}, \VUgras{linię} \textsuperscript{(d)}.

%%K t01-18-2 p
Widzimy także jedno \VUgras{działanie}, nie jest to \VUgras{dodawanie} ani
\VUgras{odejmowanie}, ani \VUgras{mnożenie}, jest to \VUgras{dzielenie}
\textsuperscript{(e)}.

%%K t01-19-1 p
19. --- Na tablicy systemu metrycznego widzimy: \VUgras{metr} \textsuperscript{(a)},
\VUgras{wagę} \textsuperscript{(b)} i jej \VUgras{odważniki}, \VUgras{litr}
\textsuperscript{(e)}, \VUgras{dekalitr} \textsuperscript{(f)}, \VUgras{decylitr} i
\VUgras{metr sześcienny} \textsuperscript{(d)}.

%%K t01-19-2 p
Uczniowie uczą się także \VUgras{nauk przyrodniczych} \textsuperscript{(11)}, zoologii
\textsuperscript{(a)}, botaniki \textsuperscript{(b)}, geologii \textsuperscript{(c)} i
\VUgras{historii} \textsuperscript{(12)}.

%%K t01-19-3 p
W szafie są przyrządy do \VUgras{fizyki} \textsuperscript{(15)} i do \VUgras{chemii}
\textsuperscript{(16)}.

%%K t01-19-4 p
Na szafie są: \VUgras{kula} \textsuperscript{(14)}, \VUgras{stożek} i
\VUgras{globus} \textsuperscript{(13)}.

%%K t01-19-5 p
Nad tablicami wisi \VUgras{mapa} przedstawiająca pięć części świata:
\VUgras{Europę} \textsuperscript{(g)}, \VUgras{Azję} \textsuperscript{(e)},
\VUgras{Afrykę} \textsuperscript{(f)}, \VUgras{Amerykę} \textsuperscript{(ab)} i
\VUgras{Oceanię} \textsuperscript{(h)} oraz dwa wielkie oceany, \VUgras{Spokojny}
\textsuperscript{(c)} i \VUgras{Atlantycki} \textsuperscript{(d)}.
\VUsaut{6.0mm}
\VUfilet{22mm}
